% GENERATED FILE. Edit canonical lessons, then run npm run generate:books.
% canonical-source-hash: fnv1a64:91406236ce05b3f9
% canonical-lessons: GE-C488-anstecken, GE-C488-infektion, GE-C488-impfen, GE-C488-erholen, GE-C488-genesen

\chapter{Wieder gesund}
\label{ch:ge-wieder-gesund}
\hlchaptermodality{488}

\begin{chapteropening}
\textbf{By the end of this chapter:} \emph{I can talk about infections and getting better.}

Complete the last lesson of chapter 488: I can talk about infections and getting better.
\end{chapteropening}

\section[anstecken]{anstecken — to infect, to pass on (an illness)}

\label{lesson:GE-C488-anstecken}



\begin{quote}
Before the new one: say the German for to treat, then the German for to prescribe.
\end{quote}

\subsection*{You'll want to know: anstecken}
\textbf{anstecken} — \textquotedblleft{}to infect, to pass on (an illness)\textquotedblright{}.

\begin{quote}
Mein Bruder hat mich mit seiner Grippe angesteckt.
\end{quote}

\begin{quote}
\emph{My brother gave me his flu.}
\end{quote}

Separable: \emph{er steckt mich an}. \emph{Sich anstecken} — to catch something. Also to pin on: \emph{eine Brosche anstecken}.

\begin{cousinweb}
\emph{anstecken} is built on \emph{an} \textquotedblleft{}on\textquotedblright{} and \emph{stecken} \textquotedblleft{}to put, to stick\textquotedblright{}, inherited from Old High German \emph{steckōn} (EtymDB).
\end{cousinweb}

\begin{grammarlens}[title={What you've built}]
One more word for telling, explaining and describing at B1.
\end{grammarlens}

\subsection*{Your turn}
\begin{itemize}
\raggedright
  \item \emph{Say it:} \emph{anstecken}
  \item \emph{Say it:} \emph{anstecken}, once more
  \item \emph{Say it:} a sentence of your own with \emph{anstecken}
  \item \emph{Recall:} say the German for to treat, then the German for to prescribe, then say \emph{anstecken} again
\end{itemize}

\begin{tcolorbox}[breakable,colback=teal!4,colframe=teal!35!black,title={Before you move on}]
What does anstecken mean? (\textquotedblleft{}To infect, to pass on (an illness)\textquotedblright{}.) Say it once more.
\end{tcolorbox}
\section[die Infektion]{die Infektion — infection}

\label{lesson:GE-C488-infektion}



\begin{quote}
Before the new one: say the German for to prescribe, then the German for to infect.
\end{quote}

\subsection*{You'll want to know: die Infektion}
\textbf{die Infektion} — \textquotedblleft{}infection\textquotedblright{}.

\begin{quote}
Gegen die Infektion bekam sie ein Antibiotikum.
\end{quote}

\begin{quote}
\emph{She got an antibiotic for the infection.}
\end{quote}

Plural: \emph{die Infektionen}.

\begin{cousinweb}
\emph{Infektion} is a learned word built on \emph{infectus}, the past participle of Latin \emph{inficere} \textquotedblleft{}to dip in, to stain, to taint\textquotedblright{} (Lewis \& Short).
\end{cousinweb}

\begin{grammarlens}[title={What you've built}]
One more word for telling, explaining and describing at B1.
\end{grammarlens}

\subsection*{Your turn}
\begin{itemize}
\raggedright
  \item \emph{Say it:} \emph{die Infektion}
  \item \emph{Say it:} \emph{die Infektion}, once more
  \item \emph{Say it:} a sentence of your own with \emph{die Infektion}
  \item \emph{Recall:} say the German for to prescribe, then the German for to infect, then say \emph{die Infektion} again
\end{itemize}

\begin{tcolorbox}[breakable,colback=teal!4,colframe=teal!35!black,title={Before you move on}]
What does die Infektion mean? (\textquotedblleft{}Infection\textquotedblright{}.) Say it once more.
\end{tcolorbox}
\section[impfen]{impfen — to vaccinate}

\label{lesson:GE-C488-impfen}



\begin{quote}
Before the new one: say the German for to infect, then the German for infection.
\end{quote}

\subsection*{You'll want to know: impfen}
\textbf{impfen} — \textquotedblleft{}to vaccinate\textquotedblright{}.

\begin{quote}
Vor der Reise lassen wir uns gegen Grippe impfen.
\end{quote}

\begin{quote}
\emph{Before the trip we are getting vaccinated against flu.}
\end{quote}

\emph{Sich impfen lassen} — to get vaccinated. The noun: \emph{die Impfung}.

\begin{cousinweb}
\emph{impfen} is inherited from Old High German \emph{impfōn}, \emph{impitōn} \textquotedblleft{}to graft a plant\textquotedblright{} (EtymDB), usually traced to an early borrowing from a Late Latin word for grafting; English \emph{imp}, first \textquotedblleft{}a graft\textquotedblright{}, is a cousin.
\end{cousinweb}

\begin{grammarlens}[title={What you've built}]
One more word for telling, explaining and describing at B1.
\end{grammarlens}

\subsection*{Your turn}
\begin{itemize}
\raggedright
  \item \emph{Say it:} \emph{impfen}
  \item \emph{Say it:} \emph{impfen}, once more
  \item \emph{Say it:} a sentence of your own with \emph{impfen}
  \item \emph{Recall:} say the German for to infect, then the German for infection, then say \emph{impfen} again
\end{itemize}

\begin{tcolorbox}[breakable,colback=teal!4,colframe=teal!35!black,title={Before you move on}]
What does impfen mean? (\textquotedblleft{}To vaccinate\textquotedblright{}.) Say it once more.
\end{tcolorbox}
\section[erholen]{erholen — to get some rest; to recover (sich erholen)}

\label{lesson:GE-C488-erholen}



\begin{quote}
Before the new one: say the German for infection, then the German for to vaccinate.
\end{quote}

\subsection*{You'll want to know: erholen}
\textbf{erholen} — \textquotedblleft{}to get some rest; to recover (sich erholen)\textquotedblright{}.

\begin{quote}
Im Urlaub habe ich mich gut erholt.
\end{quote}

\begin{quote}
\emph{I had a good rest on holiday.}
\end{quote}

Reflexive: \emph{sich erholen von} — to recover from.

\begin{cousinweb}
\emph{erholen} is built on \emph{er-} and \emph{holen} \textquotedblleft{}to fetch\textquotedblright{}, which is inherited from Middle High German \emph{holen} (EtymDB): to fetch your strength back.
\end{cousinweb}

\begin{grammarlens}[title={What you've built}]
One more word for telling, explaining and describing at B1.
\end{grammarlens}

\subsection*{Your turn}
\begin{itemize}
\raggedright
  \item \emph{Say it:} \emph{erholen}
  \item \emph{Say it:} \emph{erholen}, once more
  \item \emph{Say it:} a sentence of your own with \emph{erholen}
  \item \emph{Recall:} say the German for infection, then the German for to vaccinate, then say \emph{erholen} again
\end{itemize}

\begin{tcolorbox}[breakable,colback=teal!4,colframe=teal!35!black,title={Before you move on}]
What does erholen mean? (\textquotedblleft{}To get some rest; to recover (sich erholen)\textquotedblright{}.) Say it once more.
\end{tcolorbox}
\section[genesen]{genesen — to recover, to get well}

\label{lesson:GE-C488-genesen}



\begin{quote}
Before the new one: say the German for to vaccinate, then the German for to get some rest.
\end{quote}

\subsection*{You'll want to know: genesen}
\textbf{genesen} — \textquotedblleft{}to recover, to get well\textquotedblright{}.

\begin{quote}
Nach zwei Wochen war der Patient wieder genesen.
\end{quote}

\begin{quote}
\emph{After two weeks the patient had recovered.}
\end{quote}

Formal. Perfekt with \emph{sein}: \emph{Er ist genesen.}

\begin{cousinweb}
\emph{genesen} is inherited from Old High German \emph{ginesan}, from Proto-Germanic \emph{ganesaną}, built on \emph{nesaną}, which goes back to a root meaning \textquotedblleft{}to return, to heal\textquotedblright{} (EtymDB).
\end{cousinweb}

\begin{grammarlens}[title={What you've built}]
One more word for telling, explaining and describing at B1.
\end{grammarlens}

\subsection*{Your turn}
\begin{itemize}
\raggedright
  \item \emph{Say it:} \emph{genesen}
  \item \emph{Say it:} \emph{genesen}, once more
  \item \emph{Say it:} a sentence of your own with \emph{genesen}
  \item \emph{Recall:} say the German for to vaccinate, then the German for to get some rest, then say \emph{genesen} again
\end{itemize}

\begin{tcolorbox}[breakable,colback=teal!4,colframe=teal!35!black,title={Before you move on}]
What does genesen mean? (\textquotedblleft{}To recover, to get well\textquotedblright{}.) Say it once more.
\end{tcolorbox}
